\section{Introduction and the starting conjecture}\label{sec:intro}

This paper reports what holds, after a long audit, in a research programme built on one conjecture:

\begin{quote}
\emph{Matter, light and a tachyon sector are one object, seen from three frames.}
\end{quote}

The conjecture is the starting point, not a result.  From it the author developed, over several months, a large body of mathematics
--- the \emph{Quantum Sphaera Companion}~\cite{companion}, now more than \(1300\) pages in its research version.  Its central object is a
biquaternion kernel of three complex coordinates, together with a projector, the Heisenberg compensator, that splits every state into a
symmetric and an antisymmetric part.  On this kernel the Companion builds a cascade of layers through the Hopf fibrations and the
octonions, a master equation, and a list of physical readouts: gauge groups, masses and mixing, the cosmological constant, black holes.

Since October 2026 the Companion has been audited in numbered research rounds, and every test with a decision criterion has
been pre-registered: hypothesis, data and criterion were committed before the computation (Appendix~\ref{app:method}).  The result
is mixed, and the purpose of this paper is to state it plainly.  The algebraic core is sound.  A good part of the physical superstructure
is withdrawn or holds only under named choices.  And the framework has no dynamics yet: no action varied in all fields from which its
numbers follow.

\subsection{What this paper contains}

\begin{itemize}
\item \emph{Proved structure} (Sections~\ref{sec:kernel}, \ref{sec:octonions}).  The natural norm of the kernel and its correct signature;
the Lorentz group acting on the kernels as \(SO(3,\mathbb C)\); the symmetric triple as the centred null configuration; matter- and
tachyon-type kernels as one orbit of the duality rotation, with light as its fix-sector; the compensator and its matrix form on the
electromagnetic doublet; triality; the Hopf cascade with its corrections; and the compensator matrix, with the hypotheses that a
spectral statement about it really needs.
\item \emph{A conditional construction} (Section~\ref{sec:generation}).  One generation of the Standard Model, with right-handed neutrino,
from \(\mathbb C\otimes\HH\otimes\OO\), where chirality comes from a projector of the same form as the compensator and anomaly
cancellation alone fixes it up to a frame change.  Much of this is known from the work of Dixon, Furey and others; we say what is
specific to the present framework.
\item \emph{One input and its best reading} (Section~\ref{sec:beta}).  The frame offset \(\beta_{\rm sym}\approx0.549\), on which all numerical
readouts depend, is put in.  It can be read as a ratio of two horizon countings, but this reading selects the counting that the
literature considers less well founded.
\item \emph{Readouts and four pre-registered predictions} (Section~\ref{sec:readouts}): for \(\theta_{12}\), the \(W\) mass and twice for
\(\theta_{23}\), all pending.
\item \emph{A well-posed PDE problem} (Section~\ref{sec:op5}).  A comparison model of the probe field equations has an exact gauge symmetry,
an involutive gauged symbol at order six and four hidden sixth-order conditions; its formal integrability is open.
\item \emph{A register of what does not follow} (Section~\ref{sec:not}) and the open problems, with the requirements that a dynamics would
have to meet (Section~\ref{sec:open}).
\end{itemize}

\subsection{How to read it}

Every statement carries its status: proved, computed, conditional, a reading, open, or withdrawn, and physical statements name the frame
in which they are read.  Each section ends with a status paragraph.  References of the form ``layer v4.9.0-r63'' point to the research
rounds of the Companion, where the computations, scripts and pre-registration notes are archived; the repository is public.  The
Companion records its objects in a coded language, QAM; this paper uses ordinary notation and gives the key statements in both forms in
Appendix~\ref{app:qam}.

Nothing here claims a unified theory.  The paper is meant as a clean entry point: what a reader can check, what a reader can use, and
where the work is stuck.
